\documentclass[12pt,oneside]{article}

%%%%%%%%%%%%%%%%%%%%%%%%%%%%
%%   Zusaetzliche Pakete  %%
%%%%%%%%%%%%%%%%%%%%%%%%%%%%
\usepackage{enumerate}  
\usepackage{fancyhdr}
\usepackage{a4wide}
\usepackage{graphicx}
\usepackage{palatino}
\usepackage{multirow}
\usepackage{booktabs}
\usepackage{titlesec}
\usepackage{enumitem}% http://ctan.org/pkg/enumitem

%folgende Zeile auskommentieren für englische Arbeiten
%\usepackage[ngerman]{babel}
%folgende Zeile auskommentieren für deutsche Arbeiten
\usepackage[english]{babel}

\usepackage[T1]{fontenc}
\usepackage{paratype}\renewcommand{\familydefault}{\sfdefault}
\usepackage[utf8]{inputenc}
\usepackage[bookmarks]{hyperref}
\usepackage[justification=centering]{caption}
\usepackage[style=authoryear,natbib=true,backend=biber,maxbibnames=20]{biblatex}
\usepackage{csquotes}
\bibliography{literatur}

\setlength{\parindent}{0em} 
\setlist[itemize]{noitemsep, topsep=0pt}
\setlist[enumerate]{noitemsep, topsep=0pt}


%%%%%%%%%%%%%%%%%%%%%%%%%%%%%%
%% Definition der Kopfzeile %%
%%%%%%%%%%%%%%%%%%%%%%%%%%%%%%

\pagestyle{fancy}
\fancyhf{}
\cfoot{\thepage}
\setlength{\headheight}{16pt}

%%%%%%%%%%%%%%%%%%%%%%%%%%%%%%%%%%%%%%%%%%%%%%%%%%%%%
%%  Definition des Deckblattes und der Titelseite  %%
%%%%%%%%%%%%%%%%%%%%%%%%%%%%%%%%%%%%%%%%%%%%%%%%%%%%%

\newcommand{\JMUTitle}[9]{

  \thispagestyle{empty}
  \vspace*{\stretch{1}}
  {\parindent0cm
  \rule{\linewidth}{.7ex}
  }
  \begin{flushright}
    \sffamily\bfseries\Huge
    #1\\
    \vspace*{\stretch{1}}
    \sffamily\bfseries\large
    #2\\
    \vspace*{\stretch{1}}
    \sffamily\bfseries\small
    #3
  \end{flushright}
  \rule{\linewidth}{.7ex}

  \vspace*{\stretch{1}}
  \begin{center}
    \includegraphics[width=2in]{siegel} \\
    \vspace*{\stretch{1}}
    \Large #5 \\

    \vspace*{\stretch{2}}
   \large Lehrstuhl f\"{u}r Wirtschaftsinformatik und Business Analytics\\
    \large Universität Würzburg\\
    \vspace*{\stretch{1}}
    \large Betreuer:  #8 \\[1mm]
    \large Assistent:  #9 \\[1mm]
    \vspace*{\stretch{1}}
    \large #6, den #7
  \end{center}
}

\titlespacing*{\section}
{0pt}{3.5ex plus 1ex minus .2ex}{.2ex plus .2ex}
\titlespacing*{\subsection}
{0pt}{1.5ex plus 1ex minus .2ex}{.2ex plus .2ex}
\titlespacing*{\subsubsection}
{0pt}{1.5ex plus 1ex minus .2ex}{.2ex plus .2ex}

%%%%%%%%%%%%%%%%%%%%%%%%%%%%
%%  Beginn des Dokuments  %%
%%%%%%%%%%%%%%%%%%%%%%%%%%%%

\begin{document}
  \JMUTitle
      {Time Series Anomaly Detection Benchmarking}        % Titel der Arbeit
      {Philip Spaier}                        % Vor- und Nachname des Autors
      {3110375}
      
      {Seminararbeit } % Art der Arbeit
      {W"urzburg}                           % Ort der Erstellung
      {05.04.2025}                          % Tag der Abgabe
      {Prof. Dr. Gunther Gust}           % Name des Erstgutachters
      {Viet Nguyen} % Name des/der betreuenden Assistent/Assistentin
      
  \clearpage

\lhead{}
\pagenumbering{Roman} 
    \setcounter{page}{1}

\tableofcontents
\clearpage

\addcontentsline{toc}{section}{\listfigurename}
\listoffigures

\addcontentsline{toc}{section}{\listtablename}
\listoftables
\clearpage

\setlength{\parskip}{0.5em} 


%%%%%%%%%%%%%%%%%%%%%%%%%%%%
%%  Kurzzusammenfassung   %%
%%%%%%%%%%%%%%%%%%%%%%%%%%%%
\lhead{Abstract}
\section*{Abstract}

Eine Kurzzusammenfassung der Vorgehensweise und der wesentlichen Ergebnisse.

Allgemeine Merkmale
\begin{itemize}
    \item Objektivität: Es soll sich jeder persönlichen Wertung enthalten.
    \item Kürze: Es soll so kurz wie möglich sein.
    \item Verständlichkeit: Es weist eine klare, nachvollziehbare Sprache und Struktur auf.
    \item Vollständigkeit: Alle wesentlichen Sachverhalte sollen enthalten sein.
    \item Genauigkeit: Es soll genau die Inhalte und die Meinung der Originalarbeit wiedergeben.
\end{itemize}{}


%%%%%%%%%%%%%%%%%%%%%%%%%%%%
%%  Einstellungen  %%
%%%%%%%%%%%%%%%%%%%%%%%%%%%%
\clearpage
\pagenumbering{arabic}  
    \setcounter{page}{1}
\lhead{\nouppercase{\leftmark}}

%%%%%%%%%%%%%%%%%%%%%%%%%%%%
%%  Hauptteil  %%
%%%%%%%%%%%%%%%%%%%%%%%%%%%%


\section{Literature Review} \label{einleitung}
Time Series Anomaly Detection, as a subcategory of the broader field of Anomaly Detection, has seen increased attention since the start of the twenty first century. With the internet having established itself as a persistent and omnipresent force in every imaginable aspect of human life, time series data can be found in abundance. Modern developments in Internet-of-Things (IoT) applications, the digitization of financial data, and a massive rise in the consumption of streaming services have contributed to an exponential growth of time series data [source needed]. This in turn has made the manual search of potential anomalies in many fields completely infeasible, leading to an increased demand for automated anomaly detection methods. While there is a continuously growing repertoire of such automated detection methods, the lack of a generally accepted and reliable benchmark makes not just further developments but also the selection of appropriate models difficult. In the following sections of this literature review, I will provide the reader with a better understanding of context independent Time Series Anomaly Detection, the most commonly applied methods, and the current state of benchmarking.


\subsection{Time Series Data and Anomaly Detection Definition}
Time Series Data, as used in the rest of this thesis, shall be defined as follows: a sequence of data or observations, typically indexed by or associated with specific timestamps, collected in chronological order over a period of time.
For the purpose of analysis, continuous signals must be converted into individual data points. Each datapoint can either represent a binary state (1 or 0), be a numerical value measured on a ratio scale (eg. number of occurrences), or a numerical value measured on an interval scale (eg. temperature on a Celsius scale).
A time series with a dimensionality of one (only a single feature) will be referred to as "univariate", while a time series with higher dimensionality (multiple features) will be referred to as "multivariate". \\
An anomaly will be defined as follows: an abnormal, rarely occurring data point or sequence, that has to be be detectable with exclusively context independent methods. Individual anomalous data points will be referred to as "point based" anomalies. Multiple consecutive anomalous points, each of which might be unremarkable on their own, while displaying unusual behavior as a sequence,  will be referred to as "sequence based" or "collective" anomalies \parencites[p.~3]{liu2024elephant}[p.~8]{chalapathy2019deeplearninganomalydetection}. A separate category of anomalies would be context dependent ones. Those are data points or sequences, possibly indistinguishable from normal ones if analyzed without context, but if combined with additional information about the field or time series, are considered anomalous \citep[pp.~7-8]{chalapathy2019deeplearninganomalydetection}. Context dependent anomalies will not be topic of the research presented here. \\
Given those definitions, Time Series Anomaly Detection is therefore the task of correctly and autonomously identifying anomalies within a given time series.

Insert Images of point vs sequence.

\subsection{Relevant Fields}

The following is an overview of fields relying on Time Series Anomaly Detection. It is a non-exhaustive list, simply highlighting some of the most prominent use cases to provide context. \

\textbf{Illicit Activity and Fraud Detection:} With the global financial system relying primarily on digital transactions, it has become crucial to detect fraudulent activities as quickly and accurately as possible. A particularly obvious example is credit card fraud, creating an estimated yearly loss in the billions of dollar \parencites[p.~2]{FinFraud}. Companies like Visa and Mastercard put great emphasis on being able to detect anomalous transactions in real time to then analyses them and prevent potential harm to their customers \citep{VisaFr}. While credit card fraud is a prominent application, the scope of financial anomaly detection extends significantly further, playing a critical role in the operations of stock exchanges, brokerage firms, and banks. These institutions leverage anomaly detection techniques to identify various illicit activities, ensure market integrity, manage operational risks, and comply with stringent regulatory requirements \citep{DBSurv}. \

\textbf{Healthcare:} Healthcare critically relies on analyzing physiological signals, such as those captured by the electrocardiogram (ECG), which provides vital time series data reflecting the heart's electrical activity. While historically, ECG analysis has focused on identifying established patterns of known heart diseases, this approach often fails to detect rare or atypical anomalies that do not fit predefined categories, potentially missing critical conditions. To address this issue, Time Series Anomaly Detection has been introduced for the purpose of detecting such rare anomalies that would go unnoticed by conventional pattern classification \parencites[p.~1-2]{jiang2024anomalydetectionelectrocardiogramsadvancing}.

\textbf{Website Traffic:} A common threat faced by web-services are so called Denial or Service (DoS) and Distributed Denial of Service (DDoS) attacks. These include hitting a webserver with so many requests that the systems becomes inoperational and can no longer service legitimate users \citep{BSI1}. A significant challenge in detecting these attacks is that the malicious traffic can often mimic normal network traffic, making it difficult for traditional packet-based intrusion detection systems or statistical methods reliant on fixed thresholds to accurately identify attacks, especially when they are hidden within legitimate flows. Time series analysis allows systems to observe and distinguish the instant changes in network traffic that indicate an attack, even when individual packets or simple statistics are insufficient. Time series anomaly detection provides a means to autonomously identify and localize potentially harmful deviations within the network traffic and thereby ensure the availability and reliability of services \parencites[pp.~1-2]{FOULADI2020102587}.

\subsection{Detection Methods}
\subsection{Performance Metrics}

\subsection{State of Benchmarking}

\section{Dataset Analysis}

\textbf{Das ist fett gedruckter Text}.

\textit{Das ist kursiver Text}.


Auflistungen sind oft hilfreich für die Strukturierung:
\begin{itemize}
    \item Erster Eintrag
    \item Zweiter Eintrag
\end{itemize}

Nummerierte Aufzählungen sind oft hilfreich für Reihenfolgen:
\begin{enumerate}
    \item Erster Eintrag
    \item Zweiter Eintrag
\end{enumerate}

\section{Replication of TSB-AD Benchmark Results}

\section{Dataset Creation}

\section{Conclusion}

\section{Zitieren und Referenzieren}

Beiträge in Fachzeitschriften wie \citet{clemen1989combining} oder Konferenzartikel wie \citet{he2017mask} werden auf diese Weise im Text zitiert. In anderen Fällen möchte man aber in Klammern zitieren \citep{clemen1989combining}, auch mit mehreren Autoren \citep{clemen1989combining,baumol1958warehouse,he2017mask}.

Bei Monographien muss eine Seitenzahl mit angegeben werden \citep[S. 28]{chollet2018deep}.

So wird eine Webquelle zitiert: \citet{shiny1}. Es kann bei kurzen Informationen im Internet aber auch reichen die Adresse\footnote{\url{https://shiny.rstudio.com/tutorial/written-tutorial/lesson1/}} als Fußnote einzubetten.

So werden andere Teile der Arbeit referenziert: Kapitel \ref{einleitung}, Gleichung \ref{eq:1} zeigt...

So verweisen wir auf eine Fußnote \footnote{dies ist eine Fußnote}.

\section{Abbildungen}

Abbildungen erfordern das package \textit{graphicx}. 
Idealerweise verwendet man Vektorgrafiken oder hochaufgelöste Bitmaps. 
Eine gute Variante ist das Verwenden von PDFs.

\begin{figure}[h]
    \centering
    \includegraphics[width=0.3\textwidth]{siegel.pdf}
    \caption{Siegel der Universität}
    \label{fig:my_label}
\end{figure}


\section{Tabellen}

Die Tabular-Umgebung gibt die Anzahl Spalten an, deren Orientierung, Breite und evtl. Zwischenlinien. 


\begin{table}[ht]
    \centering
    \caption{Meine Tabelle}
        \begin{tabular}{ cccc } 
        \toprule
        col1 & col2 & col3 \\
        \midrule
        \multirow{3}{4em}{Multiple row} & cell2 & cell3 \\ 
        & cell5 & cell6 \\ 
        & cell8 & cell9 \\ 
        \bottomrule
    \end{tabular}
    \label{tab:countries}
\end{table}

\section{Formeln}

\begin{equation}
    \sum_{i=1}^N x_i
    \label{eq:1}
\end{equation}



%%%%%%%%%%%%%%%%%%%%%%%%%%%%
%% Literaturverzeichnis wird 
%% automatisch eingefügt
%%%%%%%%%%%%%%%%%%%%%%%%%%%%
\clearpage
\lhead{}
\printbibliography
\addcontentsline{toc}{section}{\bibname}


%%%%%%%%%%%%%%%%%%%%%%%%%%%%
%% Anhang (optional) 
%%%%%%%%%%%%%%%%%%%%%%%%%%%%
\clearpage
\appendix
\section{Anhang A}

%%%%%%%%%%%%%%%%%%%%%%%%%%%%
%% Eidesstattliche Erklärung
%% muss angepasst werden 
%% in Erklaerung.tex
%%%%%%%%%%%%%%%%%%%%%%%%%%%%
\newpage
\begin{otherlanguage}{ngerman}
\thispagestyle{empty}
\section*{}
\thispagestyle{empty}
Hiermit versichere ich, die vorliegende Arbeit selbstständig verfasst und keine anderen als die angegebenen Quellen und Hilfsmittel benutzt sowie die Zitate deutlich kenntlich gemacht zu haben.
\newline
Ich erkläre weiterhin, dass die vorliegende Arbeit in gleicher oder ähnlicher Form noch nicht im Rahmen eines
anderen Prüfungsverfahrens eingereicht wurde.
\vspace{4\baselineskip}\\
Würzburg, den \today \hfill Philip Spaier
\vspace{4\baselineskip}\\
\end{otherlanguage}

\end{document}
